%% Elementwise proofs: no choice of a family of witnesses is used.

\section{Powerset quantifiers, pullbacks, and equality}

\begin{annexitem}{annex:p10.quant-set.1}{thm:p10-powerset-adjunctions}{Inverse images and their two adjoints}
\begin{proof}
\paragraph{The fibres and reindexing.}
Intersection, union, complement, and the two bounds obey the
Boolean laws by membership in subsets of \(X\).
For \(W,U,V\subseteq X\),
\[
 W\cap U\subseteq V
       \quad\Longleftrightarrow\quad
 W\subseteq(X\setminus U)\cup V.
\]
Thus the specified implication is the Heyting residual.
For \(X=\varnothing\) the same equivalence holds in its
one-element fibre.

For \(x\in X\), membership in \(f^*V\) is membership of \(f(x)\)
in \(V\).  This observation proves monotonicity and preservation
of unions, intersections, and complements:
\[
\begin{aligned}
 f^*\Bigl(\bigcup_j V_j\Bigr)&=\bigcup_j f^*V_j,\\
 f^*\Bigl(\bigcap_j V_j\Bigr)&=\bigcap_j f^*V_j,\\
 f^*(Y\setminus V)&=X\setminus f^*V.
\end{aligned}
\]
The empty union and intersection give preservation of
\(\varnothing\) and top.  Combining union and complement
also proves preservation of Heyting implication.
The identity and composite inverse-image laws follow by
testing respectively \(x\) and \(g(f(x))\) for membership.
Hence \(P\), with \(P(f)=f^*\), is a contravariant functor
whose maps preserve all Heyting operations.

\paragraph{The two adjunctions.}
The condition \(\exists_fU\subseteq V\) says exactly that
\(f(x)\in V\) for every \(x\in U\).
This is \(U\subseteq f^*V\), proving the left adjunction.
Similarly, \(f^*V\subseteq U\) says that whenever \(f(x)\in V\)
one has \(x\in U\).  Equivalently, for every \(y\in V\) all
elements of its fibre belong to \(U\), which says
\(V\subseteq\forall_fU\).
Both quantifier maps are monotone: enlarging \(U\) preserves
an existential witness and preserves containment of a fibre.

If two maps satisfy the left adjunction, their values at
each \(U\) have the same upper bounds \(V\); use each value
as such an upper bound to obtain both inclusions.
If two maps satisfy the right adjunction, their values
have the same lower bounds \(V\), and the same argument
proves their equality.
This proves uniqueness without any assumption on nonempty fibres.

\paragraph{Function laws and Boolean duality.}
For an identity function, every fibre at \(x\) is \(\{x\}\),
so both quantifiers are the identity.
For \(f:X\to Y\), \(g:Y\to Z\), and \(z\in Z\),
\[
\begin{aligned}
 z\in\exists_g\exists_fU
 &\Longleftrightarrow
     \exists y\in Y\ \exists x\in U\
                       (f(x)=y\ \text{and }g(y)=z)\\
 &\Longleftrightarrow
     \exists x\in U\ (g(f(x))=z).
\end{aligned}
\]
This proves the existential composition law.
For the universal law,
\[
\begin{aligned}
 z\in\forall_g\forall_fU
 &\Longleftrightarrow
   \forall y\in Y\
     \bigl(g(y)=z\Longrightarrow
       \forall x\in X\ (f(x)=y\Longrightarrow x\in U)\bigr)\\
 &\Longleftrightarrow
   \forall x\in X\ (g(f(x))=z\Longrightarrow x\in U).
\end{aligned}
\]
In the forward implication take \(y=f(x)\); in the reverse,
apply the last condition to any \(x,y\) with \(f(x)=y\)
and \(g(y)=z\).  No witnesses for empty fibres are required.

A fibre is contained in \(U\) exactly when it has no element
in \(X\setminus U\).  This proves
\(\forall_fU=Y\setminus\exists_f(X\setminus U)\).
Existential quantification commutes with unions because an
element of \(\bigcup_jU_j\) belongs to some \(U_j\).
Universal quantification commutes with intersections because
being in every \(U_j\) can be tested for every element of
the same fibre.  These statements include empty indexed
families: the corresponding values are \(\varnothing\)
and \(Y\).  The remaining formulas follow directly from
the image and fibre definitions.

If the domain is empty, inverse image always gives
\(\varnothing\), its left adjoint sends \(\varnothing\)
to \(\varnothing\), and its right adjoint sends
\(\varnothing\) to the whole codomain.  If the codomain is
empty, a function into it can exist only when its domain
is empty; all maps between the resulting one-element
fibres are then the unique maps.  Thus every stated
identity covers the empty sets as well.
\end{proof}
\end{annexitem}

\begin{annexitem}{annex:p10.quant-set.2}{thm:p10-powerset-base-change}{Fibrewise base change and Frobenius}
\begin{proof}
\paragraph{The fibre comparison in a pullback.}
Fix \(y'\in Y'\).  Commutativity shows that \(g'\) restricts to
\[
 (f')^{-1}(\{y'\})\longrightarrow f^{-1}(\{g(y')\}).
\]
This function is a bijection when the square is a pullback.
Indeed, given \(x\) with \(f(x)=g(y')\), the two constant
maps from a singleton with values \(x,y'\) form a cone
over \(f,g\).  The pullback property supplies a unique point
\(x'\in X'\) with \(g'(x')=x\) and \(f'(x')=y'\).
Existence gives surjectivity of the fibre comparison,
and uniqueness gives injectivity.
The argument concerns each compatible pair individually.

\paragraph{Existential and universal base change.}
For \(U\subseteq X\),
\[
\begin{aligned}
 y'\in g^*(\exists_fU)
 &\Longleftrightarrow
      \exists x\in U\ (f(x)=g(y'))\\
 &\Longleftrightarrow
      \exists x'\in X'\
          (f'(x')=y'\ \text{and }g'(x')\in U)\\
 &\Longleftrightarrow
      y'\in\exists_{f'}((g')^*U).
\end{aligned}
\]
The middle equivalence is supplied by the fibre bijection.
Likewise, the fibre of \(f\) over \(g(y')\) is contained in
\(U\) if and only if every element of the fibre of \(f'\)
over \(y'\) maps under \(g'\) into \(U\).
These are exactly the conditions
\[
             y'\in g^*(\forall_fU)
               \quad\Longleftrightarrow\quad
             y'\in\forall_{f'}((g')^*U).
\]
An empty fibre has no existential witness and satisfies
the universal condition.  A bijection of fibres preserves
both facts, so the proof also covers the empty cases.
Since membership agrees for each \(y'\), both operator
identities follow.  No family of witnesses is selected.

\paragraph{What commutativity alone supplies.}
For any commuting square, if \(x'\) witnesses membership
of \(y'\) in \(\exists_{f'}((g')^*U)\), then \(g'(x')\)
belongs to \(U\) and
\[
                    f(g'(x'))=g(f'(x'))=g(y').
\]
It witnesses membership in \(g^*(\exists_fU)\).
If \(y'\in g^*(\forall_fU)\), then every \(x'\) with
\(f'(x')=y'\) has \(f(g'(x'))=g(y')\), and therefore
\(g'(x')\in U\).  This proves the automatic universal
inclusion in the stated direction.
The example with \(X'=\varnothing\) evaluates the two
opposite sides by these same fibre definitions:
the fibre above its sole parameter point is empty in
\(X'\) and a singleton in \(X\).
For \(U=X\) the existential inclusion is strict; for
\(U=\varnothing\) the universal inclusion is strict.

\paragraph{Frobenius.}
For \(y\in Y\), membership in
\(\exists_f(U\cap f^*V)\) means that some \(x\in U\)
has \(f(x)=y\) and \(f(x)\in V\).
Equivalently, \(y\in V\) and some \(x\in U\) has \(f(x)=y\).
This is membership in \(\exists_fU\cap V\), proving both
inclusions.  In particular, no element in an empty fibre
can supply the witness on either side.
\end{proof}
\end{annexitem}

\begin{annexitem}{annex:p10.quant-set.3}{thm:p10-powerset-equality}{The parameterized diagonal adjunction}
\begin{proof}
The image of the top subset \(A\) along \(\Delta_A\) is
exactly the set of pairs \((a,a)\).
Consequently
\[
 i\in\langle s,t\rangle^*\delta_A
   \quad\Longleftrightarrow\quad
 (s(i),t(i))\in\delta_A
   \quad\Longleftrightarrow\quad s(i)=t(i).
\]
Taking the diagonal itself gives
\(\Delta_A^*\delta_A=A\).

For the left-associated parameter product, a point
\(((i,a),b)\) belongs to \(p^*U\cap r^*\delta_A\) exactly
when \((i,a)\in U\) and \(a=b\).
This means that it is \(e(i,a)\) for a point \((i,a)\)
of \(U\), and is exactly membership in \(\exists_eU\).
Thus
\[
                        \exists_eU=p^*U\cap r^*\delta_A.
\]
Apply the already proved adjunction
\(\exists_e\dashv e^*\) to get
\[
          p^*U\cap r^*\delta_A\subseteq V
              \quad\Longleftrightarrow\quad U\subseteq e^*V.
\]
This proves the adjunction at every predicate \(U\) and
every parameter set \(I\).
If \(I\) or \(A\) is empty, all the relevant domain and
triple-product sets are empty and the same membership
equivalences hold in their one-element powerset fibres.

It remains to verify the indicated doctrine's full package.
The category of small sets is locally small and has the
chosen finite products.  Its powerset fibres are small
Boolean, hence Heyting, algebras, and inverse images
preserve all their Heyting operations by
Theorem~\ref{thm:p10-powerset-adjunctions}.
The identity and composite inverse-image laws make these
maps a contravariant functor.  Every projection has both
adjoints, since every set map has them.

For a parameter substitution \(h:I'\to I\), the square with
top map \(h\times1_A:I'\times A\to I\times A\), bottom map
\(h\), and vertical projections is a pullback.
Explicitly, a compatible pair \(i'\in I'\), \((i,a)\in I\times A\)
has \(h(i')=i\), and its unique lifting is \((i',a)\).
For a cone of functions \(u:B\to I'\), \(v:B\to I\times A\)
with \(h\circ u=\pi_I\circ v\), the unique mediating function
is \(b\mapsto(u(b),\pi_A(v(b)))\).  Its two composites are
\(u\) and \(v\) by the displayed compatibility, proving the
pullback universal property for every \(B\).
The two Beck--Chevalley identities of
Theorem~\ref{thm:p10-powerset-base-change} therefore give the
required parameter-product identities.  Its Frobenius
formula holds for these projections as well.
Finally the preceding diagonal calculation proves exactly
the parameterized equality axiom of
Definition~\ref{def:p10-doctrine-equality}, with the stated
left-associated maps.  These verifications prove the
Boolean first-order doctrine conclusion.

The empty context uses the terminal singleton.
For the unique function \(q:\varnothing\to1\), the definitions
give \(\exists_q(\varnothing)=\varnothing\) and
\(\forall_q(\varnothing)=1\).  These values justify the
stated empty-sort distinction without imposing any
inhabitance assumption.  All arguments use membership,
unique maps, and individual witnesses; none selects a
family of witnesses by the Axiom of Choice.
\end{proof}
\end{annexitem}
